% filepath: c:\Users\danae\Documents\piraeus_vice\report_q4.tex

\documentclass[12pt]{article}
\usepackage{amsmath}
\usepackage{amssymb}
\usepackage{graphicx}
\usepackage{booktabs}
\usepackage{hyperref}
\usepackage[margin=1in]{geometry}

\title{Q4. Linear classifier (linear networks)}
\author{Dimitris Lazanas, Alexis Kastanaras, Danai Charzaka}

\begin{document}

\section{Introduction}

This report documents the implementation of a multiclass linear classifier using Stochastic Gradient Descent (SGD) with logarithmic loss and L2 regularization to predict the killer ID for criminal incidents. The model learns linear decision boundaries that separate different killer classes in the feature space.

\section{Mathematical Derivations}

\subsection{Linear Classifier Model}

A linear classifier learns a linear decision function that maps input features to predicted class scores. The fundamental model is:

\begin{equation}
f(x) = Wx + b
\end{equation}

where $W \in \mathbb{R}^{d \times k}$ is the weight matrix, $b \in \mathbb{R}^{k}$ is the bias vector, $x \in \mathbb{R}^{d}$ is the input feature vector, and $k$ is the number of classes (in this case, 8 different killer IDs).

For a given input sample, the predicted class is determined by taking the argmax over all class scores:

\begin{equation}
c = \arg\max_{k} f_k(x)
\end{equation}

This means we select the class $k$ that produces the highest score $f_k(x)$. The weight matrix $W$ is learned during training via gradient descent, with each row corresponding to one class and each column corresponding to one feature. A larger weight value indicates that feature is more important for predicting that particular class.

\subsection{One-Hot Encoding of Target Labels}

One-hot encoding converts categorical labels (killer IDs) into a numerical format suitable for regression-based training. For a dataset with 8 killer classes (K0 through K7), each label is converted to a binary vector of length 8 with exactly one element set to 1 (indicating the true class) and all others set to 0.

For example, if a sample belongs to killer K2, its one-hot encoded representation is:
\begin{equation}
y_{\text{onehot}} = [0, 0, 1, 0, 0, 0, 0, 0]^T
\end{equation}

This encoding allows us to treat multiclass classification as a regression problem where we predict a probability distribution over classes.

\subsection{Sum of Squared Errors Loss Function}

The Sum of Squared Errors (SSE) loss measures the squared difference between the true one-hot encoded labels and the predicted probabilities. For a single sample $i$, the SSE is:

\begin{equation}
L_i(W, b) = \sum_{k=1}^{K} (y_{i,k} - P(k|x_i))^2
\end{equation}

where:
\begin{itemize}
    \item $y_{i,k} \in \{0, 1\}$ is the one-hot encoded target value for class $k$ (1 if true class, 0 otherwise)
    \item $P(k|x_i)$ is the predicted probability for class $k$ obtained from \texttt{predict\_proba()}
    \item $K = 8$ is the total number of killer classes
\end{itemize}

The average SSE loss over all $n$ training samples is:

\begin{equation}
L(W, b) = \frac{1}{n} \sum_{i=1}^{n} \sum_{k=1}^{K} (y_{i,k} - P(k|x_i))^2
\end{equation}

The goal of training is to minimize this loss so that predicted probabilities closely match the one-hot encoded targets. Ideally, for a correctly predicted sample, the probability of the true class should approach 1 while probabilities for all other classes should approach 0.

\subsection{L2 Regularization}

L2 regularization (also called ridge regularization) prevents overfitting by penalizing large weight values. The regularization term is:

\begin{equation}
R(W) = \alpha \sum_{j=1}^{d} \sum_{k=1}^{K} W_{j,k}^2
\end{equation}

where $\alpha$ is the regularization strength parameter. When $\alpha$ is large, the model is heavily penalized for having large weights, forcing it to learn simpler decision boundaries. When $\alpha$ is small, the model is allowed to have larger weights and thus learn more complex boundaries. The relationship between $\alpha$ and the inverse parameter $C$ is:

\begin{equation}
\alpha = \frac{1}{C}
\end{equation}

\subsection{Total Loss Function with Regularization}

The complete objective function that combines SSE loss and L2 regularization is:

\begin{equation}
J(W, b) = \frac{1}{n} \sum_{i=1}^{n} \sum_{k=1}^{K} (y_{i,k} - P(k|x_i))^2 + \alpha \sum_{j=1}^{d} \sum_{k=1}^{K} W_{j,k}^2
\end{equation}

During training, we minimize this objective function using stochastic gradient descent (SGD). The first term encourages the model to fit the training data, while the second term prevents overfitting by constraining the magnitude of weights.

\subsection{Optimization via Stochastic Gradient Descent}

Stochastic Gradient Descent iteratively updates the model parameters in the direction that reduces the loss:

\begin{equation}
W^{(t+1)} = W^{(t)} - \eta \frac{\partial J}{\partial W}
\end{equation}

\begin{equation}
b^{(t+1)} = b^{(t)} - \eta \frac{\partial J}{\partial b}
\end{equation}

where $\eta$ is the learning rate that controls how large each update step is, and $t$ denotes the iteration number. The model performs updates for up to \texttt{max\_iter=1000} iterations until convergence or reaching the maximum number of iterations.

\section{Methodology}

\subsection{Data Preprocessing}

\subsubsection{Train/Validation Split}

The raw dataset is split into TRAIN and VAL (validation) sets using a pre-existing ``split'' column in the data. This allows us to evaluate hyperparameter choices on held-out data without seeing the final test set. The TRAIN set is used for both model training and hyperparameter tuning via cross-validation, while the VAL set provides an unbiased estimate of final model performance.

\begin{equation}
X_{\text{train}} = \{x_i : \text{split}_i = \text{``TRAIN''}\}, \quad X_{\text{val}} = \{x_i : \text{split}_i = \text{``VAL''}\}
\end{equation}

\subsubsection{Feature Selection and Encoding}

The following columns are removed from the feature matrix as they are not useful for prediction:
\begin{itemize}
    \item \texttt{incident\_id}: Unique identifier for each incident (no predictive value)
    \item \texttt{split}: Train/validation indicator (only for splitting, not for prediction)
    \item \texttt{killer\_id}: Target variable (used to create $y$, not part of input $X$)
    \item \texttt{weapon\_code}, \texttt{scene\_type}, \texttt{weather}: Categorical features already one-hot encoded in previous preprocessing steps
\end{itemize}

The remaining features consist of numeric features (normalized via MinMaxScaler to [0,1] range) and one-hot encoded categorical features. These form the input feature matrix $X \in \mathbb{R}^{n \times d}$.

\subsubsection{Target Variable Preparation}

The killer ID labels are converted to one-hot encoding using \texttt{pd.get\_dummies()}, creating a binary matrix $y_{\text{onehot}} \in \{0, 1\}^{n \times 8}$ where each row corresponds to a sample and each column to a killer class.

\subsection{Model Architecture}

The model uses a two-layer architecture for multiclass classification:

\subsubsection{Base Estimator: SGDClassifier}

\texttt{SGDClassifier} implements linear classification with various loss functions and regularization options:
\begin{itemize}
    \item \textbf{Loss Function}: \texttt{loss='log\_loss'} uses logarithmic loss (also called logistic loss), which is appropriate for probabilistic classification. This loss function outputs well-calibrated probabilities via \texttt{predict\_proba()}.
    \item \textbf{Penalty}: \texttt{penalty='l2'} applies L2 regularization to the weight matrix
    \item \textbf{Regularization Strength}: \texttt{alpha} parameter controls the strength of L2 regularization
    \item \textbf{Optimization}: Uses stochastic gradient descent for efficient training on large datasets
    \item \textbf{Max Iterations}: \texttt{max\_iter=1000} limits training to 1000 epochs
    \item \textbf{Random State}: \texttt{random\_state=42} ensures reproducibility
\end{itemize}

\subsubsection{Multiclass Wrapper: OneVsRestClassifier}

Since \texttt{SGDClassifier} is inherently a binary classifier, we wrap it with \texttt{OneVsRestClassifier} to handle the 8-class multiclass problem. This wrapper trains one binary classifier per class, where each classifier learns to distinguish one killer class from all others. During prediction, all 8 classifiers produce scores, and the final class is determined by taking the argmax.

\subsection{Hyperparameter Tuning with GridSearchCV}

\subsubsection{Parameter Search Space}

We search over the regularization strength parameter $\alpha$ (equivalently, the inverse parameter $C$):

\begin{equation}
C \in \{0.001, 0.01, 0.1, 1, 10, 100, 1000, 10000, 100000\}
\end{equation}

\begin{equation}
\alpha = \frac{1}{C} \in \{1000, 100, 10, 1, 0.1, 0.01, 0.001, 0.0001, 0.00001\}
\end{equation}

This gives us 9 different regularization strength levels to evaluate.

\subsubsection{Cross-Validation Strategy}

We use 5-fold cross-validation on the TRAIN set to evaluate each $\alpha$ value. In each fold, 4/5 of the TRAIN data is used for training and 1/5 is used for validation. This process repeats 5 times with different train/validation splits, and the average performance across folds is used to evaluate each hyperparameter.

\subsubsection{Scoring Metrics}

Two metrics are computed during cross-validation:

\textbf{Primary Metric - Accuracy}:
\begin{equation}
\text{Accuracy} = \frac{\text{Number of correct predictions}}{\text{Total number of predictions}}
\end{equation}

This is the standard classification accuracy metric. We use accuracy as the primary metric for model selection (``refit='Accuracy'``), meaning we choose the $\alpha$ value that maximizes CV accuracy.

\textbf{Secondary Metric - Sum of Squared Errors}:

A custom scorer is implemented to compute SSE on the CV folds:

\begin{equation}
\text{SSE}_{\text{fold}} = \sum_{i=1}^{n_{\text{fold}}} \sum_{k=1}^{8} (y_{i,k} - P(k|x_i))^2
\end{equation}

This scorer uses \texttt{predict\_proba()} to get predicted probabilities and compares them against one-hot encoded labels. The scorer returns negative SSE (since GridSearchCV maximizes), and we monitor this to understand the loss landscape.

\subsubsection{Custom SSE Scorer Implementation}

The \texttt{sse\_scorer} function computes SSE by:
\begin{enumerate}
    \item Obtaining predicted probabilities from the estimator using \texttt{predict\_proba(x)}
    \item Converting categorical labels $y$ to one-hot encoding using \texttt{pd.get\_dummies()}
    \item Using \texttt{reindex()} to ensure all 8 killer classes are represented, even if some classes are missing in a particular CV fold
    \item Computing the squared differences and summing them
    \item Returning $-\text{SSE}$ for maximization in GridSearchCV
\end{enumerate}

\subsection{Model Training and Evaluation}

\subsubsection{Training on TRAIN Set}

After hyperparameter tuning, the best model (with the optimal $\alpha$ value) is automatically re-trained on the full TRAIN set by GridSearchCV. This final model is stored in \texttt{best\_estimator\_} and is ready for evaluation.

\subsubsection{Validation Set Evaluation}

The trained model is evaluated on the held-out VAL set using three metrics:

\textbf{Classification Accuracy}:
\begin{equation}
\text{Accuracy}_{\text{VAL}} = \frac{1}{m} \sum_{i=1}^{m} \mathbb{1}[\hat{c}_i = c_i^*]
\end{equation}

where $m$ is the number of validation samples, $\hat{c}_i$ is the predicted class, and $c_i^*$ is the true class.

\textbf{Sum of Squared Errors}:
\begin{equation}
\text{SSE}_{\text{VAL}} = \sum_{i=1}^{m} \sum_{k=1}^{8} (y_{i,k} - P(k|x_i))^2
\end{equation}

\textbf{Confusion Matrix}:
\begin{equation}
\text{CM}_{(i,j)} = \sum_{n=1}^{m} \mathbb{1}[\hat{c}_n = j \text{ and } c_n^* = i]
\end{equation}

The confusion matrix is an $8 \times 8$ array where element $(i, j)$ counts the number of samples with true class $i$ that were predicted as class $j$. Diagonal elements represent correct classifications, while off-diagonal elements represent misclassifications.

\subsection{Decision Boundary Visualization}

\subsubsection{PCA Dimensionality Reduction}

To visualize high-dimensional decision boundaries in 2D, we apply Principal Component Analysis (PCA) to reduce the feature space:

\begin{itemize}
    \item \textbf{Components}: PCA is fitted with 2 components to extract the two directions of maximum variance
    \item \textbf{Fitting Data}: PCA is fitted on the entire dataset $X$ (not just TRAIN) to obtain a consistent transformation across TRAIN and VAL splits
    \item \textbf{Explained Variance}: Each component explains a certain percentage of the total variance in the data
    \item \textbf{Projection}: Both TRAIN and VAL data are projected onto the 2D PCA space
\end{itemize}

\subsubsection{Decision Boundary Computation}

To visualize decision boundaries, a dense mesh grid is created covering the 2D PCA space:

\begin{equation}
\text{Grid} = \{(x_1, x_2) : x_1 \in [x_{\min} - 0.5, x_{\max} + 0.5], x_2 \in [y_{\min} - 0.5, y_{\max} + 0.5]\}
\end{equation}

with a step size of $h = 0.05$. For each point in this mesh grid, both models (Q4 linear and Q3 GaussianNB) make predictions:

\textbf{Q4 Linear Model Predictions}:
\begin{equation}
Z_{\text{linear}}(x_1, x_2) = \arg\max_k f_k([x_1, x_2]^T)
\end{equation}

\textbf{Q3 Gaussian Naïve Bayes Predictions}:
\begin{equation}
Z_{\text{GNB}}(x_1, x_2) = \arg\max_k P(k | [x_1, x_2]^T)
\end{equation}

\subsubsection{Visualization and Overlay}

The visualization overlays predictions from both models:

\begin{enumerate}
    \item \textbf{GaussianNB Decision Areas}: Semi-transparent filled contours (\texttt{alpha=0.5}) showing regions where GaussianNB predicts each class
    \item \textbf{Linear Decision Boundaries}: Filled contours showing regions where the linear classifier predicts each class
    \item \textbf{Validation Data Points}: Scatter plot of validation samples colored by true killer class, with black edges for visibility
\end{enumerate}

This overlay allows visual comparison of how the linear and non-linear classifiers partition the feature space. Linear models produce straight-edged boundaries, while GaussianNB produces smooth, curved boundaries.

\section{Results}

\subsection{Hyperparameter Tuning}

GridSearchCV evaluated 9 different regularization strengths using 5-fold cross-validation on the TRAIN set. The results are summarized in Table \ref{tab:hp_tuning}.

\begin{table}[h]
\centering
\begin{tabular}{lr}
\toprule
\textbf{Parameter} & \textbf{Value} \\
\midrule
Best $C$ & 10000.0 \\
Best $\alpha$ & 0.000100 \\
Best CV Accuracy & 0.9279 \\
Minimum SSE (CV) & 950.6061 \\
\bottomrule
\end{tabular}
\caption{Hyperparameter tuning results from GridSearchCV using 5-fold cross-validation on the TRAIN set. The optimal regularization parameter is very small ($\alpha = 0.0001$), indicating weak regularization. This suggests the model can benefit from learning more complex weight patterns without overfitting.}
\label{tab:hp_tuning}
\end{table}

The best $C$ value of 10000.0 (equivalently $\alpha = 0.0001$) was selected because it achieved the highest cross-validation accuracy of 92.79\%. This weak regularization allows the model to learn more complex decision boundaries while still generalizing well to unseen data.


\subsection{Validation Set Performance}

The best model is re-trained on the full TRAIN set and evaluated on the held-out VAL set. Performance metrics after running the program are shown below. Screenshots containing these metrics are

VAL Accuracy & 0.9248 \\
VAL SSE & 126.7239 \\

Final model performance on the held-out validation set using optimal hyperparameters 
($\alpha = 0.0001$).\\

The validation accuracy of 92.48\% is very close to the best CV accuracy (92.79\%), indicating minimal overfitting and good generalization.}
\label{tab:final_eval}
\end{table}

The small difference between CV accuracy (92.79\%) and VAL accuracy (92.48\%) suggests the model generalizes well without significant overfitting. The validation SSE of 126.7239 represents the total squared error across all validation samples and all 8 killer classes.

\subsection{Confusion Matrix}

The confusion matrix (Table \ref{tab:confusion}) provides detailed insights into which killer classes are correctly classified and which pairs are frequently confused.

\begin{table}[h]
\centering
\begin{tabular}{c|rrrrrrrr}
\toprule
\textbf{True/Pred} & K0 & K1 & K2 & K3 & K4 & K5 & K6 & K7 \\
\midrule
K0 & 11 & 0 & 2 & 0 & 1 & 2 & 1 & 0 \\
K1 & 1 & 61 & 0 & 0 & 0 & 0 & 0 & 0 \\
K2 & 1 & 0 & 480 & 2 & 17 & 1 & 0 & 7 \\
K3 & 0 & 0 & 0 & 32 & 2 & 0 & 8 & 2 \\
K4 & 0 & 0 & 5 & 1 & 23 & 0 & 3 & 3 \\
K5 & 0 & 0 & 1 & 0 & 0 & 132 & 0 & 0 \\
K6 & 0 & 1 & 3 & 8 & 2 & 0 & 140 & 0 \\
K7 & 0 & 0 & 2 & 2 & 5 & 0 & 0 & 7 \\
\bottomrule
\end{tabular}
\caption{Confusion matrix on validation set. Rows represent true killer IDs, columns represent predicted killer IDs. Diagonal elements (correct predictions) are emphasized. The matrix reveals which killer classes are well-separated and which pairs are frequently confused.}
\label{tab:confusion}
\end{table}

\subsubsection{Per-Class Performance Analysis}

Analyzing the diagonal of the confusion matrix reveals performance for each killer class:

\begin{itemize}
    \item \textbf{K0}: 11/17 correct (64.7\%) - Low accuracy, with misclassifications spread across K2, K4, K5, K6
    \item \textbf{K1}: 61/62 correct (98.4\%) - Excellent performance, nearly perfectly separated
    \item \textbf{K2}: 480/491 correct (97.8\%) - Best performer by absolute count, well-separated from most classes
    \item \textbf{K3}: 32/45 correct (71.1\%) - Moderate performance, with 8 misclassifications to K6
    \item \textbf{K4}: 23/48 correct (47.9\%) - Poorest performer, heavily confused with K2 (17 samples)
    \item \textbf{K5}: 132/134 correct (98.5\%) - Second best, nearly perfectly separated
    \item \textbf{K6}: 140/155 correct (90.3\%) - Good performance, with some confusion to K3 and K2
    \item \textbf{K7}: 7/19 correct (36.8\%) - Very poor performance, heavily confused with other classes
\end{itemize}

\subsubsection{Confusion Patterns}

The most significant confusion patterns are:

\begin{enumerate}
    \item \textbf{K2 and K4}: Mutual confusion with 17 samples from K4 misclassified as K2. This is the strongest confusion signal, suggesting these killer IDs have similar incident patterns in the feature space.
    \item \textbf{K3 and K6}: Secondary confusion pattern with 8 samples from K3 misclassified as K6.
    \item \textbf{Under-represented classes}: K0 and K7 have very low per-class accuracy (64.7\% and 36.8\% respectively), likely due to having only 17 and 19 samples respectively.
    \item \textbf{Well-separated classes}: K1 (61/62 correct) and K5 (132/134 correct) are nearly perfectly classified, indicating these killer IDs have very distinct incident patterns.
\end{enumerate}

\section{Key Figures}

\begin{figure}[h]
\centering
\includegraphics[width=0.9\textwidth]{decision_boundaries_q4_vs_q3.png}
\caption{Decision boundaries comparison: Q4 linear classifier (filled contours) overlaid with Q3 GaussianNB decision areas (semi-transparent background) in 2D PCA space. The first PCA component explains 71.84\% of the variance, while the second explains 16.65\%, collectively capturing 88.49\% of total variance. Validation samples are plotted as colored scatter points with black edges, revealing the spatial distribution of killer classes and the effectiveness of each model's decision boundaries.}
\label{fig:decision_boundaries}
\end{figure}

The visualization clearly shows the difference between linear and non-linear decision boundaries. The linear classifier (Q4) produces straight-edged regions separated by linear hyperplanes, while GaussianNB (Q3) produces smooth, curved boundaries that follow the distribution of each class. In areas where classes are well-separated (e.g., K1, K5), both models achieve high accuracy. In areas with overlapping classes (e.g., K2 and K4), both models struggle, but they make different errors due to their different decision boundary shapes.

\section{Discussion}

\subsection{Model Performance and Generalization}

The linear classifier achieves 92.48\% accuracy on the validation set, demonstrating strong overall performance. However, this aggregate metric masks significant per-class variation. The model excels at classifying well-represented, linearly-separable classes (K1: 98.4\%, K2: 97.8\%, K5: 98.5\%) but struggles with under-represented (K0, K7) and overlapping (K4) classes.

The small gap between CV accuracy (92.79\%) and VAL accuracy (92.48\%) is a positive sign, indicating the model generalizes well without significant overfitting. The weak regularization ($\alpha = 0.0001$) proves sufficient for this problem, suggesting the training data is of high quality with minimal noise.

\subsection{Regularization and Model Complexity}

The optimal $\alpha = 0.0001$ represents weak L2 regularization, allowing the weight matrix to develop large values where necessary. This indicates that:

\begin{enumerate}
    \item The underlying killer classification problem benefits from allowing the model to learn complex patterns
    \item The feature space is not excessively high-dimensional, so there is no need for aggressive regularization
    \item The TRAIN set is large enough and clean enough to support training such a model without overfitting
\end{enumerate}

\subsection{Linear vs. Non-Linear Boundaries}

The linear classifier assumes that classes can be separated by straight hyperplanes in the feature space. The confusion matrix and visualization reveal:

\begin{itemize}
    \item \textbf{Linear Advantages}: Interpretability (weights indicate feature importance), computational efficiency, and simplicity
    \item \textbf{Linear Limitations}: May struggle when classes form complex, non-linear patterns (as seen with K0, K4, K7)
    \item \textbf{GaussianNB Comparison}: Non-linear boundaries from GaussianNB provide additional flexibility to fit curved decision boundaries that follow class distributions
\end{itemize}

The strong performance of K1, K2, and K5 suggests these classes are approximately linearly separable. The poor performance of K0, K4, and K7 suggests these classes either lack sufficient training data or have complex non-linear decision boundaries.

\subsection{Class Imbalance Effects}

The dataset exhibits significant class imbalance:

\begin{itemize}
    \item \textbf{Well-represented}: K2 (491 samples), K6 (155 samples), K5 (134 samples)
    \item \textbf{Under-represented}: K0 (17 samples), K7 (19 samples), K4 (48 samples)
\end{itemize}

The model achieves high accuracy on well-represented classes but lower accuracy on under-represented ones. This is expected behavior: with more training samples, the model can learn more reliable decision boundaries. Under-represented classes have less data to learn from and are more susceptible to noise.

\subsection{Confusion Pattern Analysis}

The strongest confusion is between K2 and K4, with 17 K4 samples misclassified as K2. This suggests these killer IDs have very similar incident patterns in the feature space (time, location, victim age, weather, etc.). To improve separation between K2 and K4, one would need either:

\begin{enumerate}
    \item Better features that capture more distinct patterns between these killers
    \item A non-linear classifier that can create curved boundaries to separate overlapping regions
    \item Collecting more K4 samples to provide the model with better learning signals
\end{enumerate}

\section{Conclusion}

The multiclass linear classifier with logarithmic loss and L2 regularization achieves solid performance on the killer identification task:

\begin{itemize}
    \item \textbf{Validation Accuracy}: 92.48\%
    \item \textbf{Best CV Accuracy}: 92.79\%
    \item \textbf{Optimal Regularization}: $\alpha = 0.0001$ (equivalently $C = 10000$)
    \item \textbf{Validation SSE}: 126.7239
\end{itemize}

The model demonstrates strong generalization, with the small difference between CV and VAL accuracy indicating minimal overfitting. Performance is particularly strong for well-represented, linearly-separable classes (K1, K2, K5) and weaker for under-represented classes (K0, K7) and those with overlapping patterns (K4).

The linear decision boundaries are interpretable and computationally efficient, making this approach suitable for operational use. However, the ~47.9\% accuracy on K4 and ~36.8\% on K7 suggest room for improvement through either feature engineering, non-linear models, or addressing class imbalance.

Comparison with Q3 (GaussianNB) will reveal whether non-linear boundaries provide sufficient improvement to justify the added complexity, and whether a hybrid ensemble approach combining both models might achieve even better results.

\end{document}